% =========================================================
%  config/environments.tex
%  Lingkungan (environment) khusus buku ajar berbasis OBE.
%  Setiap kotak dipakai konsisten pada seluruh bab sehingga
%  pembaca mengenali fungsi tiap bagian.
% =========================================================

\tcbuselibrary{breakable,skins,theorems}

% --- Gaya dasar kotak buku ------------------------------
\tcbset{
  bookbox/.style={
    breakable,
    sharp corners,
    boxrule=0.7pt,
    left=7pt, right=7pt, top=5pt, bottom=5pt,
    fonttitle=\bfseries\small,
    coltitle=white,
    titlerule=0.4pt,
    toptitle=2pt, bottomtitle=2pt,
    before skip=9pt, after skip=9pt,
    before upper={\setlength{\parskip}{0.35em}}
  }
}

% --- Tujuan pembelajaran bab ----------------------------
\newtcolorbox{tujuanbox}[1][]{%
  bookbox, colframe=ubPrimary, colback=ubAlt,
  title={Tujuan Pembelajaran}, #1}

% --- Sub-CPMK yang dicapai bab --------------------------
\newtcolorbox{subcpmkbox}[1][]{%
  bookbox, colframe=ubAccent, colback=ubAlt,
  title={Sub-CPMK Bab Ini}, #1}

% --- Contoh dan ilustrasi -------------------------------
\newtcolorbox{contohbox}[1][]{%
  bookbox, colframe=ubGreen, colback=ubGreenBg,
  title={Contoh}, #1}

% --- Catatan penjelas -----------------------------------
\newtcolorbox{catatabox}[1][]{%
  bookbox, colframe=ubLine, colback=ubGray, coltitle=ubPrimary,
  title={Catatan}, #1}

% --- Kesalahan yang sering terjadi ----------------------
\newtcolorbox{peringatanbox}[1][]{%
  bookbox, colframe=ubRed, colback=ubRedBg,
  title={Perhatian}, #1}

% --- Tip praktik ----------------------------------------
\newtcolorbox{tipsbox}[1][]{%
  bookbox, colframe=ubWarn, colback=ubWarnBg,
  title={Tip Praktik}, #1}

% --- Lembar praktikum -----------------------------------
\newtcolorbox{praktikumbox}[1][]{%
  bookbox, colframe=ubPrimary, colback=white,
  title={Praktikum}, #1}

% --- Aktivitas belajar yang direkomendasikan ------------
\newtcolorbox{aktivitasbox}[1][]{%
  bookbox, colframe=ubAccent, colback=white,
  title={Aktivitas Pembelajaran yang Direkomendasikan}, #1}

% --- Asesmen dan instrumen penilaian --------------------
\newtcolorbox{asesmenbox}[1][]{%
  bookbox, colframe=ubAccent, colback=white,
  title={Asesmen dan Instrumen Penilaian}, #1}

% --- Checklist pencapaian kompetensi --------------------
\newtcolorbox{checklistbox}[1][]{%
  bookbox, colframe=ubGreen, colback=white,
  title={Checklist Pencapaian Kompetensi}, #1}

% --- Refleksi mahasiswa ---------------------------------
\newtcolorbox{refleksibox}[1][]{%
  bookbox, colframe=ubWarn, colback=white,
  title={Refleksi}, #1}

% --- Rangkuman bab --------------------------------------
\newtcolorbox{rangkumanbox}[1][]{%
  bookbox, colframe=ubPrimary, colback=ubAlt,
  title={Rangkuman}, #1}

% --- Teorema matematis (amsthm) -------------------------
\theoremstyle{definition}
\newtheorem{definisi}{Definisi}[chapter]
\newtheorem{contoh}{Contoh}[chapter]
\newtheorem{prosedur}{Prosedur}[chapter]

\theoremstyle{plain}
\newtheorem{teorema}{Teorema}[chapter]
\newtheorem{lemma}{Lemma}[chapter]

% --- Pemanggil judul bagian tanpa nomor ----------------
% \bagian{Tujuan Pembelajaran}
\newcommand{\bagian}[1]{%
  \needspace{3\baselineskip}%
  \section*{#1}%
  \addcontentsline{toc}{section}{#1}%
}
